\documentclass[12pt]{scrreprt}
\usepackage[utf8]{inputenc}
\usepackage{blindtext, xcolor}
\usepackage{comment}
\usepackage{enumerate}
\usepackage{booktabs}
\usepackage{multirow}
\usepackage[shortlabels]{enumitem}
\usepackage{graphicx}
\usepackage{tabularx}
\renewcommand\tabularxcolumn[1]{m{#1}}
\usepackage{mathtools}
\usepackage{mathpazo}
\usepackage{mdframed} 
\usepackage{float}
\usepackage{mdwlist}
\usepackage{alphabeta}
\usepackage{makecell}
\usepackage{pdflscape}
\usepackage{geometry}
\usepackage{colortbl}
%---------------------------------------------------------------------------
%--------------------------------------------------------------------------- 
% Full justification with typewriter font 
%--------------------------------------------------------------------------- 
\usepackage{everysel} 
\EverySelectfont{% 
\fontdimen2\font=0.4em% interword space 
\fontdimen3\font=0.2em% interword stretch 
\fontdimen4\font=0.1em% interword shrink 
\fontdimen7\font=0.1em% extra space 
\hyphenchar\font=`\-% to allow hyphenation 
} 
\usepackage[pagewise]{lineno}
%\usepackage[htt]{hyphenat} % Trennung von Typewriter Fonts
\def\linenumberfont{\normalfont\small}
\usepackage{svg}
%\usepackage{glossaries}
\usepackage[acronym,numberedsection]{glossaries}
\usepackage[nohyperlinks, printonlyused, withpage]{acronym}
\usepackage{listings}
\usepackage{color}
% deutsche Silbentrennung
%\usepackage[ngerman]{babel}
\usepackage[english]{babel}
% fuer Stichwortverzeichnis 
\usepackage{makeidx}
\usepackage{marvosym}
\usepackage{scrlayer-scrpage}
%\usepackage[headsepline,footsepline,plainfootsepline,markcase=upper]{scrlayer-scrpage}
\usepackage{url}
\usepackage{setspace}
\usepackage{csquotes}

%---------------------------------------------------------------------------
%---------------------------------------------------------------------------
% custom packages go here
\usepackage{xfrac}
%\usepackage{array}
%\usepackage{ragged2e}
\usepackage{boldline} 
\usepackage{array}
\newcolumntype{?}{!{\vrule width 2pt}}
\usepackage{tocloft}
\newcommand{\listequationsname}{}
\newlistof{myequations}{equ}{\listequationsname}
\newcommand{\myequations}[1]{%
\addcontentsline{equ}{myequations}{\protect\numberline{\theequation}#1}\par}
\newcommand\tab[1][3.5cm]{\hspace*{#1}}
%--------------------------------------------------------------------------- 

\usepackage{everysel} 
\onehalfspacing
%adi
\setlength{\headheight}{1.1\baselineskip}
%\usepackage[style=apa, backend=biber, language=ngerman]{biblatex}
\usepackage[style=apa, backend=biber, language=english]{biblatex}
%\DeclareLanguageMapping{ngerman}{ngerman-apa} 
\DeclareLanguageMapping{american}{american-apa}
\addbibresource{References.bib} 

\bibliography{References}
\renewcommand*{\chapterheadstartvskip}{\vspace*{0cm}}
% http://tex.stackexchange.com/questions/19738/why-doesnt-pagestyleempty-work-on-the-first-page-of-a-chapter
\renewcommand*\chapterpagestyle{scrheadings}
 
%\restylefloat{table}

\definecolor{dkgreen}{rgb}{0,0.6,0}
\definecolor{mauve}{rgb}{0.58,0,0.82}
\definecolor{gray}{rgb}{0.4,0.4,0.4}
\definecolor{darkblue}{rgb}{0.0,0.0,0.6}
\definecolor{cyan}{rgb}{0.0,0.6,0.6}

\lstset{frame=tb,
	captionpos=b,
	language=Java,
	aboveskip=3mm,
	belowskip=3mm,
	showstringspaces=false,
	columns=flexible,
	basicstyle={\small\ttfamily},
	numbers=none,
	numberstyle=\tiny\color{gray},
	keywordstyle=\color{blue},
	commentstyle=\color{dkgreen},
	stringstyle=\color{mauve},
	breaklines=true,
	breakatwhitespace=true,
	tabsize=3
}

\lstdefinelanguage{XML}
{
	morestring=[b]",
	morestring=[s]{>}{<},
	morecomment=[s]{<?}{?>},
	stringstyle=\color{black},
	identifierstyle=\color{darkblue},
	keywordstyle=\color{cyan},
	morekeywords={xmlns,version,type}% list your attributes here
}

% Umbenennung einer Einträge
%\renewcommand\tablename{Tabelle}
%\renewcommand{\figurename}{Abbildung}
%With babel (and English as language):
%\addto\captionsenglish{\renewcommand{\figurename}{Fig.}}
%\renewcommand\bibname{\section{Literaturverzeichnis}}
\renewcommand\bibname{\section{Bibliography}}
%\renewcommand*\contentsname{Inhaltsverzeichnis}
%\renewcommand{\glossaryname}{Sachwortverzeichnis}
%\renewcommand{\indexname}{\section{Sachwortverzeichnis}}
%\renewcommand{\acronymname}{\section{Abkürzungsverzeichnis}}
%\renewcommand{\listfigurename}{\section{Abbildungsverzeichnis}}
%\renewcommand{\listtablename}{\section{Tabellenverzeichnis}}
%\renewcommand\lstlistingname{Listing}
%\renewcommand\lstlistlistingname{\section{Listingverzeichnis}}

% Damit auch subsubsections nummeriert werden und im ToC auftauchen
\setcounter{secnumdepth}{3}
\setcounter{tocdepth}{2}

\makeindex
\pagestyle{scrheadings}
\clearscrheadfoot
\ihead{\includegraphics[height=1.69cm]{images/FH-burgenland-logo_white.png}}
\ohead{
	\\ \vspace{15px} \textnormal{\footnotesize{Department Informationstechnologie und Informationsmanagement}}
	%\par\nobreak\vspace{-8px}\makebox{\rule{\textwidth}{0.4pt}}
	\par\nobreak\vspace{-12px}\line(30,0){325}
	}
\ofoot{
	\pagemark
	\par\nobreak\vspace{-12px}\makebox[\linewidth]{\rule{\textwidth}{0.4pt}}
	\\
	\textnormal{%Bachelorstudiengang Information, Medien & Kommunikation
		%Bachelorstudiengang IT Infrastruktur-Management
		%Masterstudiengang Angewandtes Wissensmanagement
		%Masterstudiengang Business Process Engineering & Management
		Masterstudiengang Cloud Computing Engineering
		%Masterstudiengang Information Medien Kommunikation
		}
	%\vspace{-60px}
}
%\setheadtopline{1pt}
%\setheadsepline{0.4pt}
%\setfootsepline{0.4pt}
%\setfootbotline{1pt}
\setlength{\headsep}{1.0in}

%\input{kopf-und-fusszeile2}
\begin{document}
% FOLGENDE ZEILE KANN EINEN FEHLER PRODIZIEREN!
% Die Zeile soll bewirken, dass auf der Titelseite keine
% Seitennummer angegeben wird. Falls das so nicht funktioniert,
% muss ein anderer Workaround herhalten oder die Fußziele in
% zwei Varianten angelegt werden.
\pagenumbering{none}
\begin{titlepage}
	\thispagestyle{scrheadings}
	%Default Layout, clear footer.
	\ofoot{}
%\vspace*{1cm}
%\begin{flushright}
%    \includegraphics{images/FH-burgenland-logo.png}
%\end{flushright}
\noindent
%\vspace*{0.15cm}\text{Fachhochschule Burgenland GmbH}
%\\
%\vspace*{0.15cm}\text{Campus 1}
%\\
%\text{A-7000 Eisenstadt}
\vspace*{1cm}


\begin{center}  % Diplomarbeit ODER Magisterarbeit ODER Dissertation
\usefont{T1}{phv}{b}{n}
	\huge{Use case Smart-Factory with AWS IoT core }

    \vspace{3cm}

    \large{
    	IoT Frameworks\\
        IOTFR
          }
         
        \large{	~\newline \newline
        \\Master Program Cloud Computing Engineering \\
%        \\Masterstudiengang Cloud Computing Engineering \\
        SS2023
        }
  
\end{center}
\vspace{1cm}
%\vfill
	\usefont{T1}{phv}{m}{n}

  \noindent\begin{tabular}{@{}ll}
Authors:
& Christian Dragschitz (\lowercase{2210781021@fh-burgenland.at})
\\ & Andreas Gruber (\lowercase{2210781009@fh-burgenland.at})
\\ & Dominik Hasiwar (\lowercase{2210781002@fh-burgenland.at})
\\ & Mathias Rangger (\lowercase{2110781020@fh-burgenland.at})
\\ Date: & \today
\\ Version: & \docversion
\end{tabular}

\end{titlepage}

\pagenumbering{Roman}
\setcounter{page}{2}
\newpage
\tableofcontents
\clearpage

\pagenumbering{arabic}
\setcounter{page}{1}
\chapter{Einleitung}

\section{Assignment}
Setzen Sie den zugewiesenen Use Case um.
Finden Sie dafür ein geeignetes IoT Framework und begründen Sie dessen Eignung für den jeweiligen Case. In unserem Fall \ac{aws}.\\
\\
Inhalt der Seminararbeit (max. 10 Seiten):
\begin{itemize}
    \item Beschreiben Sie das IoT-Framework, Ihr Lösungskonzept und die gewählte technische Umsetzung
    \item Gehen Sie dabei auf Herausforderungen aus technischer und konzeptioneller Sicht sowie auf Security ein
\end{itemize}
Abgabe: 31.5.2023 via Moodle (Seminararbeit + Präsentation). Präsentation: 2.6.2023 – Demonstration des Use Cases erwünscht!

%Zitat \autocite[S. 20]{diekmann2010}



%	\begin{comment}

%\chapter{Verzeichnisse}
\cleardoublepage
%\appendix
%	\end{comment}
%\chapter{Verzeichnisse}
%http://stackoverflow.com/questions/1243342/how-to-avoid-a-page-break-before-start-of-bibliography
%\nocite{*}
\nocite{back2012web}
\begingroup
\let\clearpage\relax
\printbibliography[heading=bibnumbered] 
\endgroup

\newpage

%\bibliography{References}
\begingroup
\let\clearpage\relax
\listoffigures
\endgroup
\newpage

\begingroup
\let\clearpage\relax
\listoftables
\endgroup
\newpage
\chapter*{Formelverzeichnis}

Nur bei vielen Formeln und Abkürzungen erforderlich bzw. wenn von Betreuungsperson gewünscht.

\begingroup
\let\clearpage\relax
\vspace{-2.5cm}
\listofmyequations
\endgroup
\newpage


\chapter*{Abkürzungen}

\begin{acronym}
    \acro{aws}[AWS]{Amazon Web Services}
    \acro{cd}[CD]{Continuous Delivery/Continuous Deployment}
    \acro{ci}[CI]{Continuous Integration}
    \acro{iac}[IaC]{Infrastructure as Code}
    \acro{sca}[SCA]{Software Composition Analyse}
    \acro{sdlc}[SDLC]{Secure Development Lifecycle}
    \acro{saas}[SaaS]{Software as a Service}
    \acro{sast}[SAST]{Static Application Security Testing}        
\end{acronym}

\newpage
\addcontentsline{toc}{chapter}{Anhang}
\chapter*{Anhang}

Im Anhang sind all jene Informationen zu finden, die den Lesefluss der Arbeit stören würden, doch für die Vollständigkeit der Arbeit notwendig sind, z.B. Datentabellen, Fragebögen, etc.

\noindent Der Anhang unterliegt nicht mehr der Gliederung sondern wird mit Anhang A, Anhang B, … gekennzeichnet.

\end{document}
